\documentclass[aspectratio=169,xcolor=dvipsnames]{beamer}

\AtBeginSection[]
{
    \begin{frame}
        \frametitle{Visão Geral}
        \tableofcontents[currentsection]
    \end{frame}
}

\usefonttheme{professionalfonts} % using non standard fonts for beamer
\usefonttheme{serif} % default family is serif
\usepackage{fontspec}
\setmainfont{XCharter}[
    UprightFont    = *-Roman,
    BoldFont       = *-Bold,
    ItalicFont     = *-Italic,
    BoldItalicFont = *-BoldItalic,
]

\usepackage{hyperref}
\usepackage{graphicx} % Allows including images
\usepackage{booktabs} % Allows the use of \toprule, \midrule and \bottomrule in tables
\input{./slides/SimplePlusTheme/beamerthemeSimplePlus.sty}
\input{./slides/SimplePlusTheme/beamerinnerthemeSimplePlus.sty}
\input{./slides/SimplePlusTheme/beamerfontthemeSimplePlus.sty}
\input{./slides/SimplePlusTheme/beamercolorthemeSimplePlus.sty}

\usepackage{listings}
\setbeamercovered{transparent}

\usepackage{biblatex}
\addbibresource{slides/reference.bib}
\renewcommand*{\bibfont}{\footnotesize}
% Tamanho dos slides de "Leituras Recomendadas": maior que a bibliografia
% completa (\bibfont acima, que rende \tiny via a redefinição de \footnotesize)
% para destacar as leituras curadas.
\newcommand{\leiturasize}{\small}

\usepackage[brazilian ]{babel}
\usepackage{csquotes}
\usepackage{pgfplots}
\pgfplotsset{compat=1.18}

\usepackage{emoji}
\usepackage{amsmath}
\usepackage[xcharter,bigdelims,vvarbb]{newtxmath}
% newtxmath's operators font requests the fontspec family `XCharter(0)' in OT1;
% without these substitutions math digits and operator names silently fall back
% to Computer Modern (LaTeX Font Warning: Font shape `OT1/XCharter(0)/m/n' undefined).
\DeclareFontFamily{OT1}{XCharter(0)}{}
\DeclareFontShape{OT1}{XCharter(0)}{m}{n}{<->ssub * XCharter-TLF/m/n}{}
\DeclareFontShape{OT1}{XCharter(0)}{m}{it}{<->ssub * XCharter-TLF/m/it}{}
\DeclareFontShape{OT1}{XCharter(0)}{b}{n}{<->ssub * XCharter-TLF/b/n}{}
\DeclareFontShape{OT1}{XCharter(0)}{bx}{n}{<->ssub * XCharter-TLF/b/n}{}

\usepackage{bm}
\usepackage[table]{xcolor}
\usepackage{xcolor}
\usepackage{algpseudocode}

\usepackage{cancel}
\usepackage{ulem}

\usetikzlibrary{shapes}
\usetikzlibrary{arrows}
\usetikzlibrary {arrows.meta}
\usetikzlibrary{calc,positioning}
\usepgfplotslibrary{fillbetween}
\usetikzlibrary{angles,quotes}
\usetikzlibrary{tikzmark}


\definecolor{PastelRed}{HTML}{FFC1CC}
\definecolor{PastelBlue}{HTML}{C2F0FF}
\definecolor{PastelBlue2}{HTML}{3182BD}

\definecolor{PastelPink}{HTML}{FFCBE1}
\definecolor{PastelGreen}{HTML}{D6E5BD}
\definecolor{PastelYellow}{HTML}{F9E1A8}
\definecolor{PastelBlue}{HTML}{BCD8EC}
\definecolor{PastelPurple}{HTML}{DCCCEC}
\definecolor{PastelOrange}{HTML}{FFDAB4}


\renewcommand{\footnotesize}{\tiny}

\newcommand{\mespor}{
  \ifcase\month
    \or janeiro
    \or fevereiro
    \or março
    \or abril
    \or maio
    \or junho
    \or julho
    \or agosto
    \or setembro
    \or outubro
    \or novembro
    \or dezembro
  \fi
}
% \usepackage{csquotes}
% \usepackage{minted}
% \usemintedstyle{xcode}
\definecolor{vscLightBackground}{HTML}{FFFFFF}
\definecolor{vscLightText}{HTML}{000000}
\definecolor{vscLightGreen}{HTML}{008000}        % For Comments AND Numbers
\definecolor{vscLightRed}{HTML}{A31515}          % For Strings
\definecolor{vscLightTeal}{HTML}{267F99}
\definecolor{vscBlue}{HTML}{0000FF}% For main keywords (if, for, def...)
\definecolor{vscLightPytorch}{HTML}{800080}      % For PyTorch keywords (dark magenta)

\lstdefinestyle{vscode-light}{
    language=Python,
    backgroundcolor=\color{vscLightBackground},
    basicstyle=\ttfamily\scriptsize\color{vscLightText},
    keywordstyle=\color{vscLightTeal},
    commentstyle=\color{vscLightGreen},
    stringstyle=\color{vscLightRed},
    % Style for PyTorch keywords (using keyword group 2)
    keywordstyle=[2]\color{vscLightPytorch}, 
    morekeywords={self, True, False, None},
    morekeywords=[2]{
        torch, Tensor, device, dtype, to, nn, optim, autograd, utils, data,
        Module, Linear, Conv2d, ReLU, CrossEntropyLoss, Adam, SGD,
        Dataset, DataLoader, randn, zeros, ones, arange, cat, stack,
        view, reshape, sum, mean, max, min, matmul, no_grad, backward
    },
    keywordstyle=[3]\color{vscBlue}, 
    morekeywords=[3]{
        model, loss, y, y_hat, x , optimizer
    },
    frame=single,
    tabsize=1,
    showstringspaces=false,
    breaklines=true,
    captionpos=b,
    extendedchars=true,
    numbers=left,
}

\usepackage[cache=true]{minted}
% minted v3+ (TeX Live 2024+) auto-detects the output directory.
% minted v2 (Ubuntu apt TeX Live 2023) needs it set explicitly to
% match latexmkrc's $out_dir = 'build'.
\makeatletter
\@ifpackagelater{minted}{2024/01/01}{}{%
  \def\minted@outputdir{build/}%
}
\makeatother
\setminted{
    fontsize=\scriptsize,
    style=vs,
    breaklines=true,
    numbers=left,
    numbersep=5pt,
    frame=single,
}
\providecommand{\citeauthoryear}[1]{(\citeauthor{#1},~\citeyear{#1})}

\providecommand{\thetav}{\bm{\theta}}

% vetor coluna 2D com colchetes
\providecommand{\cvec}[2]{\begin{bmatrix} #1 \\ #2 \end{bmatrix}}

\providecommand{\varvector}[1]{\mathbf{#1}}
% instance
\providecommand{\X}{\varvector{X}}
\providecommand{\mlinstance}{\varvector{X}}
\providecommand{\mlinstancei}{\varvector{X}^{(i)}}
\providecommand{\mlinstancex}[1]{\varvector{X}^{(#1)}}

% target
\providecommand{\y}{\varvector{y}}
\providecommand{\mltarget}{\varvector{y}}
\providecommand{\mltargeti}{\varvector{y}^{(i)}}
\providecommand{\mltargetx}[1]{\varvector{y}^{(#1)}}

%hats
\providecommand{\yhat}{\hat{\varvector{y}}}
\providecommand{\mltargethat}{\hat{\varvector{y}}}
\providecommand{\mltargethati}{\hat{\varvector{y}}^{(i)}}
\providecommand{\mltargethatx}[1]{\hat{\varvector{y}}^{(#1)}}

% linear reg
\providecommand{\mllinearregression}{\mltargethat=\mlinstance\bm{\theta}}

\providecommand{\tikzarrowcolor}[3]{%
  \begin{tikzpicture}[remember picture, overlay]
    % Arrow body
    \path[fill=#3] ([shift={(#1,#2)}]current page.south west) 
      rectangle ++(0.3, -0.2);
    % Arrow head
    \path[fill=#3] ([shift={(#1+0.3,#2+0.1)}]current page.south west) -- 
         ++(0, -0.4) -- 
         ++(0.2, 0.2) -- 
         cycle;
  \end{tikzpicture}%
}

\providecommand{\tikzarrow}[2]{%
    \tikzarrowcolor{#1}{#2}{Red}
}



\DeclareMathOperator*{\argmax}{argmax}
\DeclareMathOperator*{\argmin}{argmin}



%----------------------------------------------------------------------------------------
%    TITLE PAGE
%----------------------------------------------------------------------------------------

\title{Deep Learning}
\subtitle{Grafo Computacional e Programação Diferenciável}

\author{Lucas Silveira Kupssinskü}

\institute
{
    Machine Learning Theory and Applications Laboratory \\
    PUCRS
}
\date{\mespor de \the\year}

%----------------------------------------------------------------------------------------
%    LOCAL TIKZ STYLES
%----------------------------------------------------------------------------------------

\tikzset{
  cgvar/.style   ={circle, draw=black, fill=PastelBlue, minimum size=0.85cm, inner sep=1pt, font=\small},
  cgop/.style    ={circle, draw=black, fill=Apricot!50, minimum size=0.85cm, inner sep=1pt, font=\small},
  cgedge/.style  ={->, >=stealth, thick, gray!80},
  cgval/.style   ={font=\scriptsize, fill=white, inner sep=1pt},
  cggrad/.style  ={font=\scriptsize, color=red, fill=white, inner sep=1pt},
}

\begin{document}

  \begin{frame}
    \titlepage
  \end{frame}

  \begin{frame}{Visão Geral}
    \tableofcontents
  \end{frame}

  %======================================================================
  \section{Definição e exemplo simples}
  %======================================================================

  %--- Slide: motivação ---
  \begin{frame}{Por que pensar em grafo computacional?}
    \begin{itemize}
      \item \textit{Backpropagation} é a aplicação repetida da regra da cadeia
        em uma função composta.
      \item Funções complicadas (uma rede profunda inteira) podem ser
        decompostas em operações elementares.
      \item Representar essa decomposição como um \alert{grafo} torna
        explícito o caminho que cada gradiente local precisa percorrer.
      \item Mesma estrutura permite ao computador fazer a derivada
        \alert{automaticamente}, ideia central de \textit{frameworks}
        como PyTorch e JAX.
    \end{itemize}
  \end{frame}

  %--- Slide: definição de grafo computacional ---
  \begin{frame}{Definição}
    \begin{itemize}
      \item Um \alert{grafo computacional} é um DAG (\textit{directed acyclic graph})
        em que:
        \begin{itemize}
          \item nós-folha representam \textbf{variáveis} (entradas, parâmetros);
          \item nós internos representam \textbf{operações} elementares (soma,
            produto, $\exp$, ReLU, $\dots$);
          \item arestas indicam o fluxo de dados, das entradas em direção à
            saída final.
        \end{itemize}
      \item Avaliar a função é fazer um \alert{forward pass}: percorrer o grafo
        em ordem topológica.
      \item Calcular gradientes em relação a cada folha é fazer um
        \alert{backward pass}: percorrer o grafo na ordem inversa, aplicando
        a regra da cadeia.
    \end{itemize}
  \end{frame}

  %--- Slide: exemplo (a+b)*c forward + backward com overlays ---
  \begin{frame}{Exemplo simples: $f(a,b,c) = (a+b)\cdot c$\footfullcite{karpathy2016cs231n}}
    \centering
    \resizebox{!}{0.62\textheight}{%
    \begin{tikzpicture}[node distance=1.2cm and 1.8cm]
      \node[cgvar] (a) {$a$};
      \node[cgvar, below=of a] (b) {$b$};
      \node[cgvar, below=of b] (c) {$c$};

      \node[cgop, right=of b, xshift=0.4cm, yshift=0.6cm] (plus) {$+$};
      \node[cgop, right=of plus, xshift=1.2cm, yshift=-0.6cm] (times) {$\times$};

      \node[right=of times, xshift=0.3cm] (out) {$f$};

      \draw[cgedge] (a) -- (plus);
      \draw[cgedge] (b) -- (plus);
      \draw[cgedge] (plus) -- (times);
      \draw[cgedge] (c) -- (times);
      \draw[cgedge] (times) -- (out);

      \onslide<2->{
        \node[cgval, left=0.1cm of a] {$2$};
        \node[cgval, left=0.1cm of b] {$3$};
        \node[cgval, left=0.1cm of c] {$4$};
        \node[cgval, above=0.1cm of plus] {$5$};
        \node[cgval, above=0.1cm of times] {$20$};
      }

      \onslide<3->{
        \node[cggrad, above=0.1cm of out] {$1$};
      }
      \onslide<4->{
        \node[cggrad, below=0.1cm of times] {$4$};
      }
      \onslide<5->{
        \node[cggrad, below=0.1cm of c] {$5$};
      }
      \onslide<6->{
        \node[cggrad, below=0.1cm of plus] {$4$};
        \node[cggrad, above left=0.1cm and 0.1cm of a] {$4$};
        \node[cggrad, below left=0.1cm and 0.1cm of b] {$4$};
      }
    \end{tikzpicture}
    }

    \medskip
    \footnotesize
    Valores em \textbf{preto} (forward), gradientes em \textbf{vermelho} (backward).\\
    \onslide<4->{Gate $\times$ comuta valores: gradiente recebe o valor da \textit{outra} entrada. \;}
    \onslide<6->{Gate $+$ distribui o gradiente igualmente.}
  \end{frame}

  %======================================================================
  \section{Forward e backward genéricos}
  %======================================================================

  %--- Slide: nó genérico forward ---
  \begin{frame}{Um nó genérico no \textit{forward}}
    \centering
    \begin{tikzpicture}[node distance=1.0cm and 1.6cm]
      \node[cgvar] (x) {$x$};
      \node[cgvar, below=of x] (y) {$y$};
      \node[cgop, right=of x, yshift=-0.6cm] (f) {$f$};
      \node[right=of f, xshift=0.2cm] (z) {$z = f(x,y)$};

      \draw[cgedge] (x) -- (f);
      \draw[cgedge] (y) -- (f);
      \draw[cgedge] (f) -- (z);
    \end{tikzpicture}
    \begin{itemize}
      \item No \textit{forward}, cada nó recebe valores de entrada e calcula
        sua saída local.
      \item As saídas alimentam os nós seguintes na ordem topológica.
      \item No final do grafo temos o valor da função, tipicamente uma
        \textit{loss} escalar $L$.
    \end{itemize}
  \end{frame}

  %--- Slide: nó genérico backward ---
  \begin{frame}{O mesmo nó no \textit{backward}}
    \centering
    \begin{tikzpicture}[node distance=1.6cm and 2.6cm]
      \node[cgvar] (x) {$x$};
      \node[cgvar, below=of x] (y) {$y$};
      \node[cgop, right=of x, yshift=-0.8cm] (f) {$f$};
      \node[right=of f, xshift=0.6cm] (z) {$z$};

      \draw[cgedge] (x) -- (f);
      \draw[cgedge] (y) -- (f);
      \draw[cgedge] (f) -- (z);

      \node[cggrad, above=0.1cm of x] {$\dfrac{\partial L}{\partial x} = \dfrac{\partial L}{\partial z}\cdot\dfrac{\partial z}{\partial x}$};
      \node[cggrad, below=0.1cm of y] {$\dfrac{\partial L}{\partial y} = \dfrac{\partial L}{\partial z}\cdot\dfrac{\partial z}{\partial y}$};
      \node[cggrad, above=0.1cm of z] {$\dfrac{\partial L}{\partial z}$};
    \end{tikzpicture}
    \begin{itemize}
      \item Cada aresta no sentido oposto carrega o \alert{gradiente
        \textit{upstream}} $\partial L / \partial z$.
      \item O nó multiplica esse gradiente pela \alert{derivada local}
        $\partial z / \partial x$ (regra da cadeia) e envia para a aresta de
        entrada.
      \item Se uma variável aparece em múltiplos nós, os gradientes que chegam
        nela são \textbf{somados}.
    \end{itemize}
  \end{frame}

  %--- Slide: padrões de gates ---
  \begin{frame}{Padrões comuns de \textit{gates}}
    \begin{columns}
      \begin{column}{.6\linewidth}
        \begin{itemize}
          \item \textbf{Soma} ($+$): \textit{distribuidor}, copia o gradiente
            \textit{upstream} para todas as entradas.
          \item \textbf{Produto} ($\times$): \textit{comutador}, multiplica o
            gradiente pelo valor da \textit{outra} entrada.
          \item \textbf{Máximo} ($\max$): \textit{roteador}, envia o gradiente
            inteiro para a entrada que venceu, zero para as demais.
          \item \textbf{Cópia} (uma variável usada duas vezes): \textit{somador}
            no backward, soma os gradientes que chegam.
        \end{itemize}
      \end{column}
      \begin{column}{.4\linewidth}
        \centering
        \begin{tikzpicture}[node distance=0.7cm and 1.0cm, font=\scriptsize]
          \node[cgop] (p) {$+$};
          \node[left=of p, yshift=0.3cm] (pa) {$a$};
          \node[left=of p, yshift=-0.3cm] (pb) {$b$};
          \node[right=of p] (po) {};
          \draw[cgedge] (pa) -- (p);
          \draw[cgedge] (pb) -- (p);
          \draw[cgedge] (p) -- (po);
          \node[cggrad, above=0.0cm of po,xshift=-0.1cm] {$g$};
          \node[cggrad, left=0.05cm of pa] {$g$};
          \node[cggrad, left=0.05cm of pb] {$g$};

          \node[cgop, below=1.6cm of p] (m) {$\times$};
          \node[left=of m, yshift=0.3cm] (ma) {$a$};
          \node[left=of m, yshift=-0.3cm] (mb) {$b$};
          \node[right=of m] (mo) {};
          \draw[cgedge] (ma) -- (m);
          \draw[cgedge] (mb) -- (m);
          \draw[cgedge] (m) -- (mo);
          \node[cggrad, above=0.0cm of mo,xshift=-0.1cm] {$g$};
          \node[cggrad, left=0.05cm of ma] {$g\,b$};
          \node[cggrad, left=0.05cm of mb] {$g\,a$};
        \end{tikzpicture}
      \end{column}
    \end{columns}
  \end{frame}

  %======================================================================
  \section{Exemplo: regressão logística}
  %======================================================================

  %--- Slide: motivação ---
  \begin{frame}{Outro exemplo: regressão logística}
    \begin{itemize}
      \item Considere a função
        $$f(\bm\theta, \bm x) = \frac{1}{1 + \exp\!\bigl(-(\theta_0 + \theta_1 x_1 + \theta_2 x_2)\bigr)}.$$
      \item Composição de \textit{gates} elementares: combinação afim,
        $\times(-1)$, $\exp$, $+1$, $1/u$.
      \item Vamos avaliar para
        $\theta_0 = -3,\; \theta_1 = -2,\; \theta_2 = 1,\;
         x_1 = -1,\; x_2 = -2$, depois rodar o backward.
      \item No fim, observaremos que vários \textit{gates} podem ser agrupados
        em uma única operação \textit{sigmoide}.
    \end{itemize}
  \end{frame}

  %--- Slide: TikZ sigmoid forward + backward ---
  \begin{frame}{Grafo da regressão logística}
    \centering
    \resizebox{!}{0.72\textheight}{%
    \begin{tikzpicture}[node distance=0.95cm and 1.15cm, font=\small]
      \node[cgvar] (t0) {$\theta_0$};
      \node[cgvar, below=of t0] (t1) {$\theta_1$};
      \node[cgvar, below=of t1] (x1) {$x_1$};
      \node[cgvar, below=of x1] (t2) {$\theta_2$};
      \node[cgvar, below=of t2] (x2) {$x_2$};

      \node[cgop, right=of t1, xshift=0.4cm, yshift=-0.4cm] (m1) {$\times$};
      \node[cgop, right=of x1, xshift=0.4cm, yshift=-0.4cm] (m2) {$\times$};

      \node[cgop, right=of m1, xshift=0.5cm, yshift=-0.4cm] (s1) {$+$};
      \node[cgop, right=of m2, xshift=0.5cm, yshift=-0.4cm] (s2) {$+$};

      \node[cgop, right=of s2, xshift=0.4cm] (neg) {$\times{-}1$};
      \node[cgop, right=of neg, xshift=0.4cm] (exp) {$\exp$};
      \node[cgop, right=of exp, xshift=0.4cm] (add1) {$+1$};
      \node[cgop, right=of add1, xshift=0.4cm] (inv) {$1/u$};

      \node[right=of inv, xshift=0.05cm] (out) {$f$};

      \draw[cgedge] (t1) -- (m1);
      \draw[cgedge] (x1) -- (m1);
      \draw[cgedge] (t2) -- (m2);
      \draw[cgedge] (x2) -- (m2);
      \draw[cgedge] (t0) -| (s1);
      \draw[cgedge] (m1) -- (s1);
      \draw[cgedge] (s1) -- (s2);
      \draw[cgedge] (m2) -- (s2);
      \draw[cgedge] (s2) -- (neg);
      \draw[cgedge] (neg) -- (exp);
      \draw[cgedge] (exp) -- (add1);
      \draw[cgedge] (add1) -- (inv);
      \draw[cgedge] (inv) -- (out);

      \onslide<2->{
        \node[cgval, above=0.05cm of t0] {$-3$};
        \node[cgval, above=0.05cm of t1] {$-2$};
        \node[cgval, above=0.05cm of x1] {$-1$};
        \node[cgval, above=0.05cm of t2] {$1$};
        \node[cgval, above=0.05cm of x2] {$-2$};
        \node[cgval, above=0.05cm of m1] {$2$};
        \node[cgval, above=0.05cm of m2] {$-2$};
        \node[cgval, above=0.05cm of s1] {$-1$};
        \node[cgval, above=0.05cm of s2] {$-3$};
        \node[cgval, above=0.05cm of neg] {$3$};
        \node[cgval, above=0.05cm of exp] {$20.09$};
        \node[cgval, above=0.05cm of add1] {$21.09$};
        \node[cgval, above=0.05cm of inv] {$0.047$};
      }

      \onslide<3->{
        \node[cggrad, right=0.05cm of out] {$1$};
        \node[cggrad, below=0.05cm of inv] {$-0.002$};
        \node[cggrad, below=0.05cm of add1] {$-0.002$};
        \node[cggrad, below=0.05cm of exp] {$-0.045$};
        \node[cggrad, below=0.05cm of neg] {$0.045$};
        \node[cggrad, below=0.05cm of s2] {$-0.045$};
        \node[cggrad, below=0.05cm of s1] {$-0.045$};
        \node[cggrad, below=0.05cm of m1] {$-0.045$};
        \node[cggrad, below=0.05cm of m2] {$-0.045$};
      }
      \onslide<4->{
        \node[cggrad, below=0.05cm of t0] {$-0.045$};
        \node[cggrad, below=0.05cm of t1] {$0.045$};
        \node[cggrad, below=0.05cm of x1] {$0.090$};
        \node[cggrad, below=0.05cm of t2] {$0.090$};
        \node[cggrad, below=0.05cm of x2] {$-0.045$};
      }
    \end{tikzpicture}
    }

    \medskip
    \footnotesize
    Valores em \textbf{preto} (forward), gradientes em \textbf{vermelho} (backward).
  \end{frame}

  %--- Slide: agrupando em sigmoide ---
  \begin{frame}{Agrupando \textit{gates}: a sigmoide}
    \begin{columns}
      \begin{column}{.55\linewidth}
        \begin{itemize}
          \item A sequência $\times{-}1 \to \exp \to +1 \to 1/u$
            implementa $\sigma(z) = 1/(1+e^{-z})$.
          \item Sua derivada tem forma fechada:
            $$\sigma'(z) = \sigma(z)\,(1 - \sigma(z)).$$
          \item Substituir os 4 \textit{gates} por um único nó \textit{sigmoide}
            simplifica o grafo, evita instabilidades numéricas e cobra apenas
            uma multiplicação no \textit{backward}.
          \item Granularidade fina (flexível) versus grossa (eficiente) é
            uma decisão recorrente em \textit{frameworks} de \textit{deep learning}.
        \end{itemize}
      \end{column}
      \begin{column}{.45\linewidth}
        \centering
        \begin{tikzpicture}[node distance=1.0cm and 1.4cm, font=\small]
          \node[cgvar] (z) {$z$};
          \node[cgop, right=of z] (sig) {$\sigma$};
          \node[right=of sig, xshift=0.1cm] (out) {$\sigma(z)$};

          \draw[cgedge] (z) -- (sig);
          \draw[cgedge] (sig) -- (out);

          \node[cggrad, below=0.1cm of z] {$\sigma(1{-}\sigma)\cdot g$};
          \node[cggrad, above=0.1cm of out, xshift=-0.1cm] {$g$};
        \end{tikzpicture}\\
        \footnotesize \textit{Forward} entrega $\sigma(z)$;
        \textit{backward} multiplica o gradiente \textit{upstream} por
        $\sigma(1{-}\sigma)$.
      \end{column}
    \end{columns}
  \end{frame}

  %======================================================================
  \section{Programação diferenciável}
  %======================================================================

  %--- Slide: autodiff ---
  \begin{frame}{Programação diferenciável}
    \begin{itemize}
      \item Se cada operação elementar conhece sua derivada local,
        o computador pode montar o grafo \textit{automaticamente} ao executar
        o código \textit{forward}.
      \item Esse processo é chamado \alert{diferenciação automática}
        (\textit{autodiff}\footfullcite{baydin2018automatic}), distinto de:
        \begin{itemize}
          \item \textit{symbolic} differentiation (manipulação algébrica de fórmulas);
          \item \textit{numerical} differentiation (diferença finita,
            $\partial f / \partial x \approx (f(x{+}h) - f(x))/h$).
        \end{itemize}
      \item Programação que usa \textit{autodiff} no laço principal recebe o nome
        de \textit{differentiable programming}: redes neurais são apenas um
        caso particular.
    \end{itemize}
  \end{frame}

  %--- Slide: forward mode vs reverse mode ---
  \begin{frame}{\textit{Forward mode} vs. \textit{reverse mode}}
    \begin{columns}
      \begin{column}{.5\linewidth}
        \textbf{Forward mode}
        \begin{itemize}
          \item Propaga derivadas \textit{junto} com o forward, da entrada para a saída.
          \item Custo proporcional ao número de \textbf{entradas}.
          \item Eficiente quando $|\text{entradas}| \ll |\text{saídas}|$.
          \item Usado em simulações físicas e análise de sensibilidade.
        \end{itemize}
      \end{column}
      \begin{column}{.5\linewidth}
        \textbf{Reverse mode (\textit{backprop})}
        \begin{itemize}
          \item Faz o forward, depois propaga gradientes da saída para as entradas.
          \item Custo proporcional ao número de \textbf{saídas}.
          \item Eficiente quando $|\text{saídas}| \ll |\text{entradas}|$, caso típico
            de redes neurais (saída escalar, milhões de parâmetros).
          \item Padrão em \textit{frameworks} de deep learning.
        \end{itemize}
      \end{column}
    \end{columns}
  \end{frame}

  %--- Slide: autograd na prática ---
  \begin{frame}{Na prática: \textit{autograd} de PyTorch\footfullcite{paszke2019pytorch}}
    \begin{itemize}
      \item Cada \texttt{Tensor} pode ter \texttt{requires\_grad=True}.
      \item Cada operação registrada constrói nós em um \textbf{grafo dinâmico},
        reconstruído a cada \textit{forward}.
      \item \texttt{loss.backward()} faz o reverse mode: percorre o grafo a
        partir da saída e acumula gradientes em \texttt{parameter.grad}.
      \item O otimizador (\texttt{SGD}, \texttt{Adam}, $\dots$) lê esses
        \texttt{.grad} e atualiza os parâmetros.
      \item JAX, TensorFlow 2 e Flax seguem o mesmo princípio, com variações:
        grafos estáticos compilados, transformações de funções (\texttt{jax.grad},
        \texttt{jax.jit}, \texttt{jax.vmap}).
    \end{itemize}
  \end{frame}

  %--- Slide: autograd em ação ---
  \begin{frame}[fragile]{\textit{autograd} em ação}
    \begin{columns}[T]
      \begin{column}{0.55\linewidth}
        \begin{lstlisting}[style=vscode-light, basicstyle=\ttfamily\scriptsize, language=Python]
import torch

a = torch.tensor(2.0, requires_grad=True)
b = torch.tensor(3.0, requires_grad=True)
c = torch.tensor(4.0, requires_grad=True)

f = (a + b) * c          # forward
f.backward()             # reverse mode

print(a.grad, b.grad, c.grad)
# tensor(4.) tensor(4.) tensor(5.)
        \end{lstlisting}
      \end{column}
      \begin{column}{0.42\linewidth}
        \begin{itemize}
          \item Mesmo grafo do primeiro exemplo: $f = (a + b)\cdot c$ com $a{=}2, b{=}3, c{=}4$.
          \item Gradientes coincidem com a derivação manual:
            $\partial f/\partial a = c = 4$, $\partial f/\partial b = c = 4$, $\partial f/\partial c = a + b = 5$.
          \item O grafo foi montado \textbf{automaticamente} ao executar a expressão; nenhum código de derivada foi escrito.
        \end{itemize}
      \end{column}
    \end{columns}
  \end{frame}

  %======================================================================
  \section*{Referências e Leituras}
  %======================================================================

  %--- Slide: leituras ---
  \begin{frame}{Leituras Recomendadas}
    \footnotesize
    \sloppy
    \begin{itemize}
      \item Notas de aula CS231n, módulo \textit{Backpropagation, Intuitions}:
        \citeauthoryear{karpathy2016cs231n},
        \href{https://cs231n.github.io/optimization-2/}{\textit{cs231n.github.io/optimization-2}}.
      \item Survey de \textit{automatic differentiation}:
        \citeauthoryear{baydin2018automatic}.
      \item Capítulo 7 de \fullcite{prince2023understanding}.
      \item PyTorch \textit{autograd}: \citeauthoryear{paszke2019pytorch} e
        documentação oficial em
        \href{https://pytorch.org/docs/stable/notes/autograd.html}{\textit{pytorch.org/docs/stable/notes/autograd}}.
    \end{itemize}
  \end{frame}

  \begin{frame}[allowframebreaks]{Referências}
    \sloppy
    \printbibliography[heading=none]
  \end{frame}

\end{document}
